% Generated from locale/tr/content/second-order-logic/metatheory/introduction.tex; do not hand-edit.


\olfileid[tr]{sol}{met}{int}

\olsection{Giriş}

Birinci dereceden mantığın bir dizi güzel özelliği vardır. Karar verilemez
olduğunu biliyoruz, ama en azından aksiyomlaştırılabilir. Başka bir
deyişle, birinci dereceden mantık için sağlam ve tam ispat sistemleri vardır.
Bu sistemler bir $\Proves$ !!a{derivability} bağıntısı doğurur ve bu bağıntı,
her $\Gamma$ !!{sentence}s kümesi ve her $!A$ !!{sentence} için
$\Gamma\Entails !A$ ancak ve ancak $\Gamma\Proves !A$ özelliğini taşır.
Özellikle bu, birinci dereceden
mantığın geçerli tümcelerinin !!{computably
  enumerable} olduğu anlamına
gelir. Hesaplanabilir bir $f\colon\Nat\to\Sent[L]$ fonksiyonu vardır ve
$f$'nin görüntüsünü tam ve yalnız $\Lang{L}$ dilinin geçerli
!!{sentence}s oluşturur. Bunun nedeni, !!{derivation}s için bir numaralandırma
bulunabilmesi ve bunlardan !!{derive} sözcüğüyle belirtilen işlemin sonunda
tek bir~!!{sentence} verenlerin
daha sonra o !!{sentence} öğesine eşlenebilmesidir. İkinci dereceden
mantık, birinci dereceden mantıktan daha ifade gücü yüksek olduğundan,
geçerliliklerini yakalamak genel olarak daha karmaşıktır. Gerçekten de ikinci
dereceden mantığın yalnızca karar verilemez olmadığını, geçerliliklerinin
!!{computably
  enumerable} bile olmadığını göstereceğiz. Bu, ikinci dereceden
mantık için sağlam ve tam bir ispat sisteminin olamayacağı anlamına gelir
(gerçi sağlam fakat eksik ispat sistemleri vardır ve bunlar önemli araştırma
nesneleridir).

Birinci dereceden mantığın iki özelliği daha vardır: tıkızdır (bir
$\Gamma$ !!{sentence} kümesinin her sonlu alt kümesi sağlanabilir olduğunda
$\Gamma$ da sağlanabilir) ve Löwenheim--Skolem Teoremi onun için geçerlidir
($\Gamma$'nın sonsuz bir modeli varsa !!a{denumerable} modeli de vardır).
Bu iki sonuç da ikinci dereceden mantıkta geçerli değildir. Bunun nedeni yine
ikinci dereceden mantığın !!{domain}s söz konusu olduğunda bunların
büyüklüklerine ilişkin, birinci dereceden mantığın ifade edemediği olguları
ifade edebilmesidir.

